% I-TREA II -- interim report on development (phase 1), 17 October 2026.
% Limits: abstract <= 150 words; main text (title page to the start of References) <= 8 pages; whole PDF <= 15 pages.
% Main text: sections/{introduction,method,setup,results,related,limitations,conclusion}.tex. It uses no internal
% codes (no decision-record ids, arm codes, set or variant names); those appear only in the appendices.
% Results are development data on 17 instances (descriptive, interim). The phase-1 result is the four-seed reading
% of the final configuration (D-031, logs/d031-reading.json); the author's go at checkpoint C1 is D-033.
% Tables in generated/ and figures in figures/ are written by tools/paper_tables.py; never edit them by hand.
% Bibliography: refs.bib, legally readable works only; paywalled works are named only in Appendix~\ref{app:paywalled}.
% Build: cd paper && OMP_NUM_THREADS=1 nice -n 19 taskset -c 10-14 latexmk -pdf -interaction=nonstopmode main.tex
\documentclass[11pt]{article}
\usepackage[margin=1in]{geometry}
\usepackage[T1]{fontenc}
\usepackage[utf8]{inputenc}
\usepackage{lmodern}
\usepackage{amsmath,amssymb}
\usepackage{booktabs}
\usepackage{array}
\usepackage{graphicx}
\usepackage{xcolor}
\usepackage{placeins}
\usepackage{tikz}
\usetikzlibrary{arrows.meta,positioning}
\usepackage[numbers,sort&compress]{natbib}
\usepackage{xurl}
\usepackage[hidelinks]{hyperref}
\graphicspath{{figures/}}
% Long instance names in \texttt{} do not hyphenate; let TeX stretch lines instead of running into the margin.
\setlength{\emergencystretch}{3em}
% Let a page hold more float area, so tables and figures stay next to their text.
\renewcommand{\topfraction}{0.85}
\renewcommand{\bottomfraction}{0.5}
\renewcommand{\textfraction}{0.1}
\renewcommand{\floatpagefraction}{0.75}

% A visible marker for text that still depends on results or decisions.
\newcommand{\pending}[1]{\textcolor{gray}{[\emph{#1}]}}
% A table written by the analysis code into generated/; a visible marker while the file does not exist yet.
\newcommand{\gentable}[1]{\IfFileExists{generated/#1.tex}{\input{generated/#1}}{%
  \begin{center}\pending{table \texttt{\detokenize{generated/#1.tex}} not yet written}\end{center}}}

\title{I-TREA II: a GPU evolutionary algorithm\\ coupled with MIP solvers ---\\
interim report on development (phase~1)}
\author{Berk Orbay\\ \small\texttt{berk.orbay@tideseed.com}\thanks{Designed and carried out with AgentRA (Agent
Research Assistant; mainly Claude Opus 5.5, working through Claude Code), under the author's direction; see
Appendix~\ref{sec:ai}. Every result of this study was produced by the code, configurations and job definitions in the
project repository, \url{https://github.com/Tideseed/i-trea-ii}.}}
\date{17 October 2026}

\begin{document}
\maketitle

\begin{abstract}
On an NVIDIA DGX Spark, I-TREA~II couples a GPU evolutionary algorithm (EA) with SCIP, exchanging solutions both
ways; the EA also takes NVIDIA cuOpt's solutions. Its population keeps a capped share of infeasible solutions whose
objective beats the incumbent's. This report covers development on 17 public instances that were chosen and
tuned on (120-s runs). On the 11 instances hard for cuOpt alone, I-TREA~II's primal integral (0: optimal at once; 1: no
useful solution) was lower than with cuOpt alone (same cores and GPU) on 7 and higher on 3 (mean 0.690 to 0.550, four
seeds). On the 15 hard for SCIP alone (one core, 6--7 October), it was lower on 12 and tied on 3 (0.889 to 0.448;
cuOpt alone on these 15: 0.517). In the largest loss, the EA delays cuOpt's first solution from 6.5--12.2 to 65\,s.
Unseen instances are tested next.
\end{abstract}

% Inserted directly after the abstract (Berk 2026-10-16: "clearly and concisely explain how these findings help NVIDIA cuOpt
% and SCIP, at the almost top of the README and paper"). Sources: logs/d031-reading.json readings[0].four (table, vs_a6f,
% vs_a0_17); cuOpt alone's mean P on SCIP's hard set (0.517) and the 8 / 8 count over the 17 computed from four.table;
% the largest loss outside cuOpt's hard set other than chromaticindex: neos-3656078-kumeu 0.085 -> 0.107 (+0.022);
% doon (2 of 4 seeds) and rail01 (1 of 4) from datasets/runs.parquet, sets c1r-ysync-ld0 and c1r-ysync-ld0-s23.
% Review fix 2026-10-16: the numbers of the abstract and Table tab:users are not repeated here.
\paragraph{What this means for SCIP and cuOpt users.}
In short: on instances hard for a solver alone, there is a real chance of a large gain, but not on every instance
(Table~\ref{tab:users}). Most of the gain over SCIP alone is what cuOpt alone also reaches on the same five cores and GPU
(mean $P$ on SCIP's hard set: SCIP alone 0.889, cuOpt alone 0.517, I-TREA~II 0.448). cuOpt on the whole machine, which is stronger
(Section~\ref{sec:limits}), has not yet been compared. Over all 17 instances, I-TREA~II was better than cuOpt alone on 8
and worse on 8. On the 6 outside cuOpt's hard set it was worse on 5, by at most 0.022 in $P$ except on
\texttt{chromaticindex1024-7} (0.08 to 0.54), where the EA delays cuOpt's first solution (Section~\ref{sec:res-loss}). On
\texttt{doon} and \texttt{rail01}, cuOpt alone found a solution on 2 and 1 of 4 seeds, I-TREA~II on none. All runs
lasted 2 minutes, on instances we tuned on.

% Short instance names in the table and the main text: ca = ca__Scenario400_reserves_1, item_placement =
% item_placement_9903, doon = neos-2746589-doon, toguru = neos-4763324-toguru (Appendix app:c1, Table tab:d031).
% SCIP row, largest gains: among the instances where I-TREA II also beats cuOpt alone (neos-960392 0.584 -> 0.042,
% ca 0.555 -> 0.261; four.table); cuOpt alone on SCIP's 15: 0.517 against SCIP alone's 0.889 (-42 %).
\newcommand{\chg}[2]{\mbox{$#1\to#2$}}
\begin{table}[t]\centering\small
\caption{Change in the primal integral $P(120)$ (lower is better; 0 = optimal at once, 1 = no solution or a gap of 1
throughout) when a solver runs inside I-TREA~II, on that solver's hard set: the instances where it alone had $P(120)\ge
0.5$ (Table~\ref{tab:terms}). Seed means: four seeds for I-TREA~II and cuOpt alone, run back to back; two for SCIP
alone. 2-minute runs on one DGX Spark; $p$: Wilcoxon signed-rank. In the SCIP row, the bracket gives the change from SCIP
alone to cuOpt alone, and the largest gains are among the instances where I-TREA~II also beats cuOpt alone. ca: the unit-commitment case; full names in Appendix~\ref{app:c1}.}
\label{tab:users}
\setlength{\tabcolsep}{4pt}
\begin{tabular}{@{}>{\raggedright}p{2.4cm}>{\raggedright}p{1.35cm}>{\raggedright}p{3.3cm}>{\raggedright}p{4.45cm}%
>{\raggedright\arraybackslash}p{3.5cm}@{}}
\toprule
solver alone & hard set & change in mean $P$ & largest gains & losses \\
\midrule
SCIP (one core; 6--7 October) & 15 of 17 & $-50\,\%$ (cuOpt alone: $-42\,\%$); 12 better, 0 worse; $p=0.0005$ &
  \mbox{neos-960392} \chg{1.00}{0.04}; ca \chg{1.00}{0.26} & none \\
cuOpt (same 5 of 20 cores and GPU) & 11 of 17 & $-20\,\%$; 7 better, 3 worse; $p=0.03$ &
  \mbox{neos-960392} \chg{0.58}{0.04}; ca \chg{0.56}{0.26}; \mbox{item\_placement} \chg{0.49}{0.22} &
  doon \chg{0.86}{1.00}; toguru \chg{0.55}{0.61}; rail01 \chg{0.98}{1.00} \\
\bottomrule
\end{tabular}
\end{table}


\section{Introduction}\label{sec:intro}
Planners and schedulers who use MIP solvers rarely wait for a proof of optimality; they need a good feasible solution
within a fixed time. How soon a solver finds one is measured by the primal integral $P(T)$~\citep{berthold2013primal},
the gap to the best known value averaged over the time budget $T$, which rewards finding good solutions early.

GPUs now sit beside CPU solvers on small machines, and they suit population search: many candidate solutions can be
evaluated and repaired at once. NVIDIA cuOpt already runs a GPU population of solutions with recombination and local
search beside its own branch-and-bound~\citep{corduk2025gpuprimal,corduk2025evoframework,nvidia2026cuoptsrc}. In an
earlier study on the same machine, simply running cuOpt beside HiGHS or SCIP lowered the primal integral by 21--30\,\%
at 60\,s, while a CPU evolutionary algorithm exchanging solutions with SCIP gave no measurable gain and lost to a second
copy of SCIP~\citep{orbay2026optopt}.

I-TREA II tests the changes that study named: a population that lives on the GPU, longer budgets (300\,s and
1{,}800\,s), and instances on which the solver's own heuristics struggle. The evolutionary algorithm (EA) runs beside
an exact solver --- SCIP as the primary partner, HiGHS as a replication partner --- and the two exchange solutions in
both directions: the solver's incumbent, very good solutions from its pool, its proven bound and its root LP solution
enter the population, and the population's better incumbents enter the solver as new starting points. From the first
I-TREA~\citep{orbay2011itrea} it keeps one principle: good infeasible solutions --- infeasible solutions whose objective
is better than the best feasible one --- are kept in the population, here as a capped share (Rank~0), beside a
protected share of the best feasible solutions (Rank~1).

\paragraph{Question.} The claim under test is that a GPU-based evolutionary metaheuristic coupled with an exact solver
can beat either one alone. On hard MIP instances, does the coupled system find good feasible solutions sooner than
\emph{the solver alone} and \emph{the EA alone}? Both comparisons must succeed for the claim to hold. Running two methods
side by side can beat either one simply by being two, so we also compare the coupling with the same two methods running
without exchange, and with the solver racing cuOpt on the same GPU (Section~\ref{sec:design}).

\paragraph{Closest prior work.} None of the parts is new. Incumbents were exchanged in both directions between a CPU
heuristic and a branch-and-bound solver running in parallel two decades ago~\citep{nwana2005cooperative,
puchinger2005cooperating}. ReXi runs local-search and large-neighbourhood workers beside SCIP workers over a shared
solution pool~\citep{mexi2026rexi}; CHAP pools the feasible and infeasible candidates of GPU and CPU
heuristics~\citep{tjusila2026chap}; and inside cuOpt a GPU population with recombination, which keeps infeasible
members and protects the best feasible one, runs beside cuOpt's own CPU branch-and-bound with solutions flowing both
ways~\citep{corduk2025evoframework,nvidia2026cuoptsrc}. Reserving a share of the population for infeasible solutions
goes back at least to IDEA~\citep{singh2011thesis}. What this study puts together is a GPU population coupled with an
\emph{external, general} MIP solver through two-way exchange, with a capped share of infeasible solutions whose
objective beats the incumbent, and with the solver's proven bound removing those that can never become feasible. We did
not find this combination in our search, which is a negative search result, not a priority claim
(Section~\ref{sec:related}).

\paragraph{Status of this draft.} The hypotheses, arms, instance-pool rules, metrics and analysis were fixed before any
experiment and are summarised in Section~\ref{sec:design}. No experiment has been run; Section~\ref{sec:results} is a
placeholder. \pending{Contributions to be stated after the results; no claim is made here.}

\section{Method}\label{sec:method}
\pending{This section describes the design as specified before any experiment; parameters are the defaults of the
pre-registered starting configuration and may change in the capped tuning phase (Section~\ref{sec:design}).}

\subsection{Problem form, violation and feasibility}
We consider $\min\, c^\top x$ subject to $r^l \le Ax \le r^u$, $l \le x \le u$ and $x_j \in \mathbb{Z}$ for
$j \in I$, read from an MPS file; maximisation is handled through a sign. The population compares \emph{keys}
(lower is better). The total violation of an individual is
\[
V(x) = \sum_i \frac{\max(r^l_i - a_i, 0) + \max(a_i - r^u_i, 0)}{s_i} + \sum_j (\text{bound violation})_j
     + \sum_{j \in I} |x_j - \operatorname{round}(x_j)|, \qquad a = Ax,
\]
with a row scale $s_i$ (the row's 2-norm by default). $V$ orders infeasible individuals; it is never the feasibility
test. An individual is \emph{feasible} if every row, bound and integrality condition holds within a tolerance under the
partner solver's own rule --- relative to $\max(|a_i|, |r_i|, 1)$ for SCIP, absolute for HiGHS --- tested row by row,
never as a sum. Objective values are always computed in double precision. Every solution passed to a solver is
re-checked on the CPU in double precision first, and every solution that enters a reported metric is re-checked by an
independent exact checker.

\subsection{The three ranks}
Each generation, the candidate set (population, offspring and solutions injected from the solver) is reduced to the
population size $P$:
\begin{enumerate}
\item Exact duplicates are removed (their provenance is merged into the copy that is kept).
\item The incumbent key $k^*$ is the better of the best feasible key in the candidate set and the solver's incumbent.
\item \textbf{Rank 1 (feasible good):} the best $\lceil y \cdot \min(\#\text{feasible}, P) \rceil$ feasible candidates
  ($y = 10\,\%$) are protected; candidates with equal keys leave Rank~1.
\item \textbf{Rank 0 (infeasible good):} infeasible candidates with key below $k^*$ (by a relative margin); those beyond
  the solver's proven bound are removed. At most $\lfloor x P \rfloor$ survive ($x = 10\,\%$): over the cap, the
  non-dominated fronts in (key, $V$) are taken in order and the overflowing front is thinned by repeatedly removing the
  point nearest to another. Candidates over the cap are removed, not demoted, so $x$ bounds the number of super-optimal
  infeasible survivors.
\item \textbf{Rank 2 (rejects):} the remaining places are filled from this generation's new candidates (offspring,
  injected solutions, immigrants) by non-dominated sorting in (key, $V$) with crowding distance~\citep{deb2002nsga2};
  the previous generation's Rank~2 does not survive. This makes the protection of Rank~1 matter: with $y = 0$ the
  best feasible solutions are not kept by rule. A shortfall is filled with immigrants.
\end{enumerate}
Before any feasible solution is known, all candidates are ranked by $V$ (a feasibility phase). Parents are chosen by
tournaments that prefer Rank~1 to Rank~0 to Rank~2, then the lower key, then the lower violation.

\subsection{Variation and repair}
Crossover (probability 0.9) uses by default a mix of \emph{agreement} (weight 0.6), uniform (0.2) and constraint-block
(0.2) crossover. In agreement crossover the child keeps the genes on which both parents agree, and repair works only on
rows that contain a disagreeing gene --- related to the combination heuristics in MIP
solvers, which fix the variables on which pool solutions agree and solve the rest as a
sub-MIP~\citep{berthold2006primal}. Mutation changes an absolute
number of integer genes per child (Poisson, mean 4) by flips, steps, moves toward the LP solution, moves inside violated
rows or moves that improve the objective. Children are then repaired by a feasibility-jump local
search~\citep{luteberget2023feasibility}: a weighted violation over the rows, one variable at a time moved to a
\emph{jump value} that minimises it, and row weights that grow at local minima. Once an incumbent exists, an extra
row $c^\top x \le k^* - \delta$ joins the weighted rows, so that repair keeps an objective better than the incumbent
instead of trading it away. The GPU implementation runs one thread block per individual with incremental updates of the
row activities; a single-core CPU implementation of the same algorithm is the reference and the control.
Singleton continuous columns are decoded in closed form; other continuous variables are moved by their jump values.

\subsection{Exchange with the solver}
Solver and EA run as separate processes released together by a start barrier and share one monotonic clock. They
exchange solutions through a Redis server on the local machine or, as a fallback, through a mailbox directory in shared
memory; the choice is fixed by a pre-registered overhead measurement. Into SCIP, solutions enter only through a
heuristic plug-in that tries each posted solution; into HiGHS, through its user-solution callback. The solver posts its
incumbent, up to eight pool solutions within 1\,\% of it, its proven bound and, once, its root LP solution; the EA posts
a feasible solution whenever it beats the solver's incumbent. cuOpt never receives solutions (accepting them disables
its presolve); it is a control arm only.

\subsection{The frontier}
Beside the ranks the EA keeps a \emph{frontier}: an archive of feasible solutions within 5\,\% of the incumbent that
are non-dominated in (objective, distance from the incumbent), thinned to 32 by crowding distance. It is an output (a
set of diverse alternative solutions, reported with its size, spread and hypervolume), an optional source of
immigrants, and in one ablation a source of additional solutions posted to the solver. With several objectives the same
archive becomes their Pareto front, as in the first I-TREA.

\section{Experimental setup}\label{sec:setup}
Phase~1 is development: the method was changed and judged openly on a small set of instances, so every number in
Section~\ref{sec:results} is descriptive. Table~\ref{tab:terms} defines the terms used from here on.

\begin{table}[t]\centering\footnotesize
\caption{Terms used in this report.}\label{tab:terms}
\begin{tabular}{>{\raggedright}p{2.6cm}>{\raggedright\arraybackslash}p{13.0cm}}
\toprule
term & meaning \\
\midrule
primal integral $P$ & $P(120)$: the gap $|v-r|/\max(|v|,|r|)$ of the best checked solution value $v$ to the reference
  $r$ (1 before the first solution or if the signs differ), averaged over the 120-s budget. \\
coupled system & the method of Section~\ref{sec:method}; its $P$ is that of the best solution held by any of its
  processes. \\
cuOpt alone & cuOpt~26.08 on the coupled system's share of the machine (5 of its 20 CPU cores and the GPU), run in
  the same window. \\
SCIP alone & SCIP~10 on one CPU core, two seeds, run on 6--7 October. \\
bounded difference $d$ & per instance, on seed means: $(P_\text{alone} - P_\text{coupled}) / \max(P_\text{alone},
  P_\text{coupled}, 0.001)$, in $[-1, 1]$; positive when the coupled system has the lower $P$. \\
hard set & the instances where a solver alone has $P \ge 0.5$: 11 for cuOpt, fixed on earlier runs of cuOpt alone
  with three cores; 15 for SCIP, on the compared runs (Section~\ref{sec:limits}). \\
hint & mean $P$ of the coupled system over that of cuOpt alone, on cuOpt's hard set; present if at most 0.95. It only
  decides whether testing continues. \\
guard & mean $d$ against cuOpt alone over all 17 instances; holds if at least $-0.05$. \\
harm check & per instance, the final configuration against the one before LP decoding and the yielding wait (13
  October); flags if their $d_i$ (as $d$) is below minus the instance's seed-to-seed band (at least 0.05). \\
\bottomrule
\end{tabular}
\end{table}

\paragraph{Machine.} One NVIDIA DGX Spark: a GB10 with 10 Cortex-X925 and 10 Cortex-A725 CPU cores, a Blackwell GPU and
121\,GiB of memory shared by CPU and GPU. The software is SCIP~10 through PySCIPOpt 6.2.1, HiGHS 1.15.1, cuOpt~26.08
and CuPy 14.2.0 on CUDA~13. No other GPU process ran beside the coupled system. Runs start only below temperature
gates, and a run whose GPU clock drops is voided and repeated (Appendix~\ref{app:harness}).

\paragraph{Instances.} We used 17 public instances (Appendix~\ref{app:c1}). Nine were chosen after a screening run of
SCIP alone. Seven are MIPLIB 2017 instances~\citep{gleixner2021miplib} on which SCIP held a large gap or found no
solution; one is a PGLib-UC unit-commitment case~\citep{pglibuc2019,knueven2020uc} and one an ML4CO item-placement
instance~\citep{gasse2022ml4co}. Eight MIPLIB instances from other families were added later. All 17 were chosen by
looking at results, the method was tuned on them, and they are excluded from later test sets. References are MIPLIB's
optimal values and provisional best known values for the other two.

\paragraph{Arms.} Several processes can do better than one simply by being several. cuOpt alone therefore gets the
coupled system's whole machine share: two Cortex-X925 cores, three Cortex-A725 cores and the GPU. Both arms of each
instance and seed ran back to back in one window, under the same memory caps. SCIP alone had fewer resources: one core,
no GPU, run on 6--7 October. Only solutions that pass an independent exact checker count in $P$.

\paragraph{Statistics.} We report mean $d$ with a 95\,\% percentile bootstrap interval over instances (10{,}000
resamples)~\citep{dicicco1996bootstrap} and the Wilcoxon signed-rank test on $d$ (ties dropped). The checkpoint
result uses seeds 0--3 (SCIP alone: 0--1); development tests used one to six seeds.

\paragraph{How changes were judged.} Each change was tested against the unchanged configuration in the same window,
under a dated record with a written reading rule. It was kept only if its checks all passed: validity, the intended
mechanism, no harm per instance, and the hint. Most rules were fixed before any run was read; Appendix~\ref{app:devrecord} lists the
exceptions. A crashed run counts for its arm, scored up to the crash.

\paragraph{The checkpoint.} At the end of phase~1 the author decides whether to go on to unseen instances, reading the
hint, guard and harm check under rules fixed before any checkpoint run was read.
\FloatBarrier

\section{Results}\label{sec:results}
\textbf{Pre-registered; no confirmatory results yet.} Only the screening and the open development runs have been run. This section will report, in the order fixed in
Section~\ref{sec:design}, every pre-registered comparison with its estimate, interval and verdict, including
comparisons that are negative or inconclusive.

\subsection{Screening and the instance pool}
\pending{Screening pass, strata, the frozen pool and references.}

\subsection{Calibration and gates}
\pending{CPU/GPU interference, exchange transport overhead, the first-feasible gate against cuOpt, local-search
throughput against the CPU reference.}

\subsection{Main claim: the coupling against either method alone}
\pending{H1a (A3 vs A0) and H1b (A3 vs A5).}

\subsection{Secondary comparisons}
\pending{H2 (uncoupled portfolio, with the live A3r check), H3 (SCIP $\parallel$ cuOpt), H4 (SCIP $\parallel$ SCIP),
H5 (tuned SCIP).}

\subsection{Good infeasible solutions}
\pending{H6 (x = 10\,\% vs 0 on the EA alone) and the descriptive lineage analysis.}

\subsection{Integrity}
\pending{Checker coverage, post accounting, crashes and voided blocks per arm.}

\subsection{Exploratory}
\pending{Protected elite ($y = 0$), CPU vs GPU EA, cuOpt alone, exchange direction and other ablations, HiGHS as the
partner, 1{,}800\,s, strata, the frontier.}

\section{Related work}\label{sec:related}
Each part of I-TREA II has prior art. Our search ran on 5 October 2026, before the first experiment; works without a
legally readable full text are listed in Appendix~\ref{app:paywalled}, not cited.

\paragraph{Keeping infeasible solutions.} Deb's feasibility rules made ``feasible first'' the common default in
constrained evolutionary optimisation~\citep{deb2000efficient,coello2002survey,mezuramontes2011survey}. Several methods
keep infeasible solutions on purpose: ASCHEA steers the feasible share~\citep{hamida2000aschea}, FI-2Pop keeps a
separate infeasible population~\citep{kimbrough2008fi2pop}, and Farmani and Wright single out infeasible solutions
better than the best feasible one~\citep{farmani2003selfadaptive}. SMES keeps the best infeasible
solution~\citep{mezuramontes2005smes}, IDEA reserves a fixed share for infeasible solutions~\citep{singh2011thesis},
and I-TREA added three ranks and an interactive decision maker~\citep{orbay2011itrea}.

\paragraph{Heuristics working with exact solvers.} Two-way exchange of incumbents between a heuristic and a
branch-and-bound solver running in parallel dates back two decades~\citep{nwana2005cooperative,puchinger2005cooperating,
puchinger2005combining,gallardo2007mabnb}. MIP solvers recombine solutions from their own pool~\citep{berthold2006primal,
hendel2022alns}, racing solvers that share incumbents gain modestly~\citep{ralphs2018parallel,
carvajal2014diversification,shinano2018fiberscip}, and ReXi runs local-search workers beside SCIP workers over a shared
pool~\citep{mexi2026rexi}.

\paragraph{GPU heuristics.} GPU LP solvers~\citep{applegate2021pdlp,lu2023cupdlp}, bound
propagation~\citep{sofranac2022domainprop} and LP-free local search~\citep{luteberget2023feasibility,lin2024localmip}
are the building blocks of current GPU MIP heuristics. cuOpt combines them~\citep{corduk2025gpuprimal} and runs a GPU
population with recombination beside its own branch-and-bound~\citep{corduk2025evoframework,nvidia2026cuoptsrc}; its
source prompted two of our tests (5 and 6). CHAP pools feasible and infeasible candidates of GPU and CPU
heuristics~\citep{tjusila2026chap}; others fix and propagate from LP solutions~\citep{kempke2025fixprop} or
search many subproblems at once~\citep{pauphilet2026pangea,sim2026ptilp}.

\paragraph{Position.} The two closest works couple differently. cuOpt's population keeps infeasible members but ranks
them behind the best feasible one, and works only with cuOpt's own branch-and-bound. ReXi shares a pool with SCIP on the
CPU and keeps no infeasible solutions. Our search found no GPU population with a capped share of infeasible solutions
better than the incumbent, coupled to an external MIP solver by two-way exchange. This is a negative search result, not
a priority claim.
\FloatBarrier

% Review fixes 2026-10-16: each caveat has one home. The crashes, the two harm flags and day-to-day drift are stated with
% their numbers in Section sec:res-c1; the combined configuration and its loosened rule in sec:res-dev. Concentration
% recomputed on the four-seed result (logs/d031-reading.json readings[0].four.table, fixed 11): neos-960392 carries
% 0.542 of the summed P reduction 1.535 (35 %); the ratio without it is 0.858. crna: D-031's note and the 2026-10-14
% diagnosis note (at two seeds about +-0.04). sing44 0.471 and item_placement 0.488: four.table p_a6f.
\section{Limitations}\label{sec:limits}
The results describe the development instances; they do not show how the method does elsewhere.
\begin{itemize}
\item \textbf{Development data.} The 17 instances were chosen and tuned on (Section~\ref{sec:res-dev}: on eight new
  ones an earlier version lost to cuOpt alone); the result rests on four seeds (SCIP alone: two).
\item \textbf{The gain over SCIP alone is not attributed.} The coupled system contains a full cuOpt run, and the
  checkpoint has no SCIP--cuOpt pair without the EA.
\item \textbf{Selection.} cuOpt's hard set was fixed on earlier runs of cuOpt alone on three cores
  (\texttt{item\_placement}, 0.488, and \texttt{sing44}, 0.471, now sit just below 0.5); SCIP's, on the compared runs.
\item \textbf{Concentration.} \texttt{neos-960392} carries 35\,\% of the summed $P$ reduction on cuOpt's hard set
  (the ratio without it: 0.858). \texttt{neos-2657525-crna} is bimodal ($\pm 0.04$ on the ratio at two seeds).
\item \textbf{Checker tolerance.} Our exact checker accepted some \texttt{glass4} solutions better than MIPLIB's proven
  optimum (one of them in the 13 October run the harm check compares against) and rejected 15 cuOpt solutions in 5 of
  the 30 calibration runs. A single acceptance rule for all arms is pending; its effect on $P$ and on the 1.31 ratio
  below is not yet known.
\item \textbf{One machine, a small share.} On five cores and the GPU, cuOpt alone's primal integral is 1.31 times
  that of GAMS's published cuOpt~26.08 run on a whole GB10 (MIPfeas
  scoring~\citep{bussieck2026mipfeas,mittelmann2026mipfeas}, 30 instances; Appendix~\ref{app:harness}), so our
  cuOpt-alone baseline is weaker than cuOpt on the whole GB10.
\item \textbf{Not measured:} HiGHS as a partner.
\end{itemize}
\FloatBarrier

% Review fix 2026-10-16: no numbers here (they are in the abstract, Table tab:users and Section sec:res-c1); the go and
% the caveats carried forward as D-033 states them.
\section{Conclusion}\label{sec:conclusion}
Phase~1 ends with the author's decision to continue: the final configuration lowered $P$ on most instances that are
hard for each solver alone, with one large loss (Section~\ref{sec:res-loss}). The two harm flags, cuOpt alone's crashes
and that loss are carried forward as stated caveats. Still open: how much of the gain needs the EA rather than cuOpt's
own solutions, and whether any of it holds on unseen instances, which the next phase tests on cuOpt~26.10 under rules
fixed before its runs.
\FloatBarrier


% References and appendices are set smaller so that the whole PDF stays within 15 pages (2026-10-16).
{\scriptsize\setlength{\bibsep}{0.5pt plus 0.3pt}
\bibliographystyle{plainnat}
\bibliography{refs}}

\appendix
\footnotesize
% Dense appendices (set here so the main text's layout is unchanged): more floats per page, tighter float,
% caption and heading spacing; run-in paragraph headings in the body size.
\setcounter{topnumber}{3}\setcounter{bottomnumber}{2}\setcounter{totalnumber}{5}
\setlength{\textfloatsep}{8pt plus 2pt minus 2pt}\setlength{\floatsep}{6pt plus 2pt minus 2pt}
\setlength{\abovecaptionskip}{3pt}\setlength{\belowcaptionskip}{2pt}
\makeatletter
\renewcommand\section{\@startsection{section}{1}{\z@}{-1.8ex \@plus -.4ex \@minus -.2ex}{0.8ex \@plus .2ex}%
  {\normalfont\large\bfseries}}
\renewcommand\paragraph{\@startsection{paragraph}{4}{\z@}{0.8ex \@plus .3ex \@minus .2ex}{-1em}%
  {\normalfont\bfseries}}
\makeatother
% Appendices: A development record, B per-instance C1 table, C measurement harness and calibration,
% D reproducibility, E AI assistance, F works without a legal free copy. They may use internal codes (D-NNN, arm
% names); the main text refers to them only by label: app:devrecord, app:c1, app:harness, sec:repro, sec:ai,
% app:paywalled. The full decision records and the dated corrections live in the repository (vault/50-decisions).
% Appendix A (rewritten 2026-10-16 by the appendix writer). Sources: the decision records vault/50-decisions/D-001 ...
% D-033; table_devrecord is generated by tools/paper_tables.py (DEVRECORD, one source comment per row). The exceptions
% list is the pre-condense design section's "Decision protocol" (fact-checked 2026-10-16) in short form.
\section{Development record}\label{app:devrecord}
Table~\ref{tab:devrecord} lists every test of phase~1, one row per decision record, numbered as in the main text.
Test~8 is the checkpoint's first reading (earlier configuration, 13 October); tests 16 and 17 are the two halves of
the result, and test~17 reads all four seeds. The framework records (D-001 to D-010,
D-013, D-019 to D-021, D-024) and the go at C1 (D-033) are, like all records with their amendments and dated
corrections, in the repository (\path{vault/50-decisions/}).

\gentable{table_devrecord}

\paragraph{Exceptions to rules fixed before the runs} (each stated in its record): tests 1, 2 and 15, in
Table~\ref{tab:changes}'s notes a--c and Section~\ref{sec:res-dev} (test~2's two deviations: a criterion built on a
premise a one-variable test refuted, and a mis-set band; test~15's thresholds were fixed with seeds 0--1 known);
test~6, adopted although one check that it worked as intended failed by its letter (the decoder's host time on one
instance, a cost that could only handicap the treatment); and test~8, whose amendment fixing the checkpoint's hint and
guard came after 7 of 58 runs had finished, none read.

% Appendix B (rewritten 2026-10-16 by the appendix writer). Source: table_d031 is generated by tools/paper_tables.py
% from logs/d031-reading.json and datasets/runs.parquet. Review fix 2026-10-16: the caveats list (D-031 Result, D-033)
% is not repeated here; each caveat has its one home in the main text (sec:res-c1, sec:limits).
\section{Per-instance checkpoint results}\label{app:c1}
Table~\ref{tab:d031} gives the phase-1 result per instance, under rules fixed before the runs (D-022, D-025, D-030,
D-031). VALID: 136 of 136 runs, cuOpt alone's three crashed runs kept as outcomes.

\gentable{table_d031}

Caveats: Sections~\ref{sec:res-c1} and~\ref{sec:limits}, and in the repository D-033.

% Appendix C (rewritten 2026-10-16 by the appendix writer, <= 1 page). Sources: the pre-condense design section's "The
% external harness and the equal footprint" and "Thermal governance as part of the measurement" (fact-checked
% 2026-10-16), D-003, D-008, D-009, D-021; the footprint table as in the C1 re-read's run records. table_calib and
% table_chromatic are generated by tools/paper_tables.py.
\section{Measurement harness and calibration}\label{app:harness}
Every compared GPU arm gets the same machine share $F$ (Table~\ref{tab:footprint}). Thermal rules are part of the
measurement, since a capped GPU would bias a comparison; the runner enforces them (D-003, D-008, D-009, D-021). The
120-s budget starts at a shared barrier, after every process has read the instance and warmed its GPU kernels.

\begin{table}[!tb]\centering\scriptsize
\caption{The machine share $F$ under the final harness. Memory caps are hard (systemd \texttt{MemoryMax} on the
host, a pool limit on the GPU); the runner's own threads run outside $F$.}\label{tab:footprint}
\setlength{\tabcolsep}{3pt}
\begin{tabular}{>{\raggedright}p{4.3cm}>{\raggedright}p{3.6cm}>{\raggedright\arraybackslash}p{7.6cm}}
\toprule
process & cores & memory caps and settings \\
\midrule
SCIP~10, one thread & Cortex-X925 core 5 & host 9\,GiB (SCIP's own limit 8\,GiB) \\
EA and its two HiGHS processes & Cortex-X925 core 6; GPU & host 6\,GiB; GPU pool 12\,GiB; CUDA yield flag \\
cuOpt~26.08, three threads & Cortex-A725 cores 1--3; GPU & host and GPU 24\,GiB each; MIP gap limits 0; thread
  variables (OMP, OpenBLAS, MKL) unset; own kernel cache \\
cuOpt alone, five threads & all five cores; GPU & as cuOpt above \\
\bottomrule
\end{tabular}
\end{table}

\paragraph{Thermal rules.} The host's fan governor caps the GPU clock at 2{,}200\,MHz after 20\,s at or above
85\,$^\circ$C at maximum fan, and requests a 30-minute rest after 30 minutes at maximum fan.
\begin{itemize}\setlength{\itemsep}{0pt}\setlength{\parskip}{0pt}
\item \textbf{Windows and start gate:} daytime windows under a lease, with the local language-model servers stopped.
  A GPU run starts below 80\,$^\circ$C (CPU-only: 88\,$^\circ$C), with no rest requested and no clock cap; rests took
  about a quarter of window time.
\item \textbf{Voids:} a block (one instance and seed, both arms) is voided, requeued once and rerun whole if the hottest zone
  reaches 92\,$^\circ$C, a GPU run's median busy SM clock is below 2{,}300\,MHz or meets a clock cap, another program's GPU
  process runs for over 2\,\% of the budget, or a kernel compiles on the clock. Below 2{,}000\,MHz (a known stuck-clock
  fault) the run set halts.
\item \textbf{Breaker:} two thermal or cap voids close the window. In the C1 re-read it fired once (15 October, D-030);
  the remaining blocks ran whole the next morning.
\end{itemize}

\paragraph{Harness measurements} (D-021), made before any run on unseen instances; they test the harness, not the
method. M4 reads outside its band, which D-021 attributes to $F$ (five threads and the GPU, against a 24-thread whole
machine) and to the machine, so the unseen-instance test's targets are anchored on our own cuOpt-alone runs.
\gentable{list_calib}
Table~\ref{tab:chromatic} gives the per-run values behind Table~\ref{tab:loss}.

\gentable{table_chromatic}

% Appendix D (rewritten 2026-10-16 by the appendix writer, <= 1 page). Sources: the previous reproducibility appendix
% (fact-checked); counts re-counted by the appendix writer on the run table of 2026-10-16 13:01 (1,497 rows, 18
% synthetic, 1,469 completed, 10 failed; 192,602 incumbents, 28 rejected in 13 runs; C1 sets 222 runs, 18,098 of
% 18,100 passed); SHA-256 of configs/frozen-v2.yaml and the job files re-computed; EA hash 835c7c5a249c re-computed
% with itrea.config.EAConfig.from_yaml; parameter values from configs/frozen-v2.yaml and src/itrea/config.py defaults.
\section{Reproducibility}\label{sec:repro}
Every number here comes from the code, configurations and job files in the project repository (to be made public
with the final version). The chain, with analysis steps pinned (\texttt{OMP\_NUM\_THREADS=1 nice -n 19 taskset -c
10-14}); each test's job file and reading command are in its decision record and the experiment index:
\begin{enumerate}\setlength{\itemsep}{0pt}\setlength{\parskip}{0pt}
\item \texttt{itrea make-jobs} (or a test's tool) writes \path{jobs/*.jsonl}: instance, arm, seed, budget,
  configuration, harness, set, block.
\item \texttt{itrea pool jobs/<file>.jsonl} writes \path{traces/raw/<set>/<run-id>/}: run record, logs, every
  incumbent.
\item \texttt{itrea build-runs} checks every incumbent and writes \path{datasets/runs.parquet} and
  \path{datasets/incumbents_checked.parquet}.
\item The C1 readings are saved by \texttt{itrea c1r-report -{}-a9 ysync-ld0 -{}-json logs/d030-reading.json}, then
  \texttt{tools/d031\_report.py -{}-d030-json logs/d030-reading.json -{}-json logs/d031-reading.json}.
\item \texttt{tools/paper\_tables.py} writes \path{paper/generated/} and \path{paper/figures/}; \texttt{-{}-check}
  compares byte for byte.
\end{enumerate}

\paragraph{Runs and the checker.} A run record holds the configuration hashes, harness, cores, caps, versions,
revision, a clean-tree flag and temperatures. The checker re-reads the model with another library, decides bounds,
integrality and objective exactly, and re-checks rows exactly where floating point could be wrong (tolerance
$10^{-6}$, all arms); only passing incumbents enter $P$. Run table of 16 October, 13:01: 1{,}479 runs (1{,}469 completed, 10 failed and scored up to the failure);
28 of 192{,}602 incumbents, in 13 runs, failed. In the four C1 run sets (222 runs, all from a clean tree), 18{,}098 of
18{,}100 passed.

\paragraph{The final configuration} (frozen at the checkpoint). \path{configs/frozen-v2.yaml}, a byte copy of
\path{configs/itrea2-dev-rc-3way-ysync-ld0.yaml} (SHA-256 \texttt{d5907fc3\ldots}): EA hash \texttt{835c7c5a249c}, A9
run hash \texttt{f345de0ab936} (the last field of every A9 run id). Run hashes of cuOpt alone and SCIP alone:
\texttt{82a371305f88}, \texttt{27f3171fac78}. A repository test fails if the frozen file carries another hash pair.
Table~\ref{tab:params} lists its parameters.

\begin{table}[!tb]\centering\scriptsize
\caption{Parameters of the final configuration: (f) set in the file, others code defaults
(\texttt{src/itrea/config.py}).}\label{tab:params}
\setlength{\tabcolsep}{3pt}
\begin{tabular}{>{\raggedright}p{1.7cm}>{\raggedright}p{6.05cm}>{\raggedright}p{1.7cm}>{\raggedright\arraybackslash}p{6.05cm}}
\toprule
part & value & part & value \\
\midrule
population & 256 (fewer if it exceeds the GPU pool); 256 children per generation & ranks & Rank~0 at most 10\,\% of
  the population; Rank~1 the best 10\,\% of the feasible \\
selection & binary tournament: Rank~1, 0, 2, then objective, then violation & variation & crossover 0.9 (agreement
  0.6, uniform 0.2, constraint block 0.2); mutation of Poisson(4) integers \\
repair & feasibility jump, at most 20{,}000 moves, stop after 1{,}000 without gain (f) & repair clock & active (f);
  horizons $2^7$--$2^{16}$ steps, or none \\
strategies & four (row scale $\times$ LP guidance) (f); redraw 0.1; each at least 5\,\% & stagnation & after 50
  generations, 25\,\% of Rank~2 replaced by immigrants \\
EA's LP & HiGHS on the EA's core (f), 60-s limit & LP stream & on (f); at most 2\,\% of SCIP's time; 16 points,
  2 per tick, 4 children each \\
LP decoding & on (f), after the EA's LP and at its priority (f); up to 4 LPs per batch, 2\,s (infeasible 1\,s, first
  10\,s) & exchange & mailbox polled every 0.25\,s; up to 8 SCIP pool solutions within 1\,\%; cuOpt's incumbents \\
GPU wait & CUDA yield flag (f) & other & frontier of 32 (output only); feasibility tolerance $10^{-6}$; GPU pool
  12\,GiB \\
\bottomrule
\end{tabular}
\end{table}

\paragraph{Re-analysis and replication.} A re-analysis must match exactly (\texttt{itrea reproduce} covers the
pre-registered chain, not yet the per-test reports or this table writer); a replication will differ (wall-clock
metrics, GPU arithmetic, thermal rests). Provisional references: the study's best on the unit-commitment instance (it
re-scores those runs) and the screening's best on the item-placement instance.

\paragraph{Environment.} Ubuntu 24.04, kernel 6.17.0-1031-nvidia, driver 580.173.02, CUDA 13.0; Python 3.12.3,
NumPy 2.4.6, SciPy 1.18.1, CuPy 14.2.0, HiGHS 1.15.1, PySCIPOpt 6.2.1 (SCIP 10.0), cuOpt 26.08, pandas 3.0.3, PyArrow
23.0.1 (pinned in \path{pyproject.toml}; each run record stores its versions).

\paragraph{Not yet public.} The repository is private until the final version. Raw runs (\path{traces/raw/}, about
3.5\,GB) and run tables are to be published as a versioned archive, as for the predecessor
study~\citep{orbay2026optopt}; step~4 regenerates the readings in \path{logs/}. Instances and GAMS's archive are
third-party data. Open (issue~\#36): the checker accepted incumbents below MIPLIB's optimum on \texttt{glass4} (11
runs, two in the C1 sets) and \texttt{mas74}; one rule for all arms is due before the second pre-registration.

\section{AI assistance statement}\label{sec:ai}
This study is designed and carried out with AgentRA (Agent Research Assistant), an AI agent working under the author's
direction. AgentRA is mainly Anthropic's Claude Opus 5.5, working through Claude Code. It drafted the design, the
pre-registration and the code, ran the literature searches (with sub-agents of the same kind), and will run the
experiments and the analysis; the author set the research question and the main claim, made or approved the decisions
recorded in the repository, and is responsible for the content.

Before any experiment, the design and the core code were reviewed by three independent reviewer agents in a Claude
Code workflow, each with one lens (algorithm, methodology, systems). All 53 findings, the response to each and the
resulting changes are recorded in the repository. \pending{External reviews of the finished study, if any, will be
named here.}

Every reference was checked for a legally readable full text; works that could not be read are listed only in
Appendix~\ref{app:paywalled}. Every number in a final version will come from the code, configurations and job
definitions in the repository and can be re-derived from the archived runs. Mistakes made during the study, including
those made by the AI agent, are recorded in the repository's incident log.

% Appendix F (written 2026-10-16 by the appendix writer; replaces the paywalled list in the old sections/appendix.tex).
% Source: vault/40-literature/Paywalled research.md § 1 (all 13 paywalled-only works of sweeps A-E), with its "cite
% instead" routes where that work is already in this report's bibliography (no new bibliography entries here).
% Two entries carry the vault's open flags: Runarsson & Yao 2000 ("needs Berk": archived author page) and Lodi &
% Tramontani 2013 (open-access status unverified). The CEC 2006 copy with an institutional-licence stamp is never linked.
\section{Works without a legal free copy}\label{app:paywalled}
Table~\ref{tab:paywalled} lists relevant works behind paywalls, with no legal free copy found. \emph{We have not read
them in full}; what we say rests on abstracts, metadata and the open works cited instead.

\begin{table}[!htb]\centering\scriptsize
\caption{Works we could not read in full, with the open work cited instead (--: none in this report).}\label{tab:paywalled}
\setlength{\tabcolsep}{3pt}
\begin{tabular}{>{\raggedright}p{9.2cm}>{\raggedright}p{4.6cm}>{\raggedright\arraybackslash}p{2.0cm}}
\toprule
work & DOI & cited instead \\
\midrule
Runarsson and Yao (2000), stochastic ranking, IEEE TEVC &
  \nolinkurl{10.1109/4235.873238} & \citep{mezuramontes2011survey}$^{a}$ \\
Ray et al.\ (2009), IDEA, book chapter & \nolinkurl{10.1007/978-3-642-00619-7_7} & \citep{singh2011thesis} \\
Singh et al.\ (2008), IDEA, AI 2008 & \nolinkurl{10.1007/978-3-540-89378-3_11} & \citep{singh2011thesis} \\
Isaacs et al.\ (2008), CHEA, IDEA's precursor, IEEE CEC & \nolinkurl{10.1109/CEC.2008.4631171} &
  \citep{singh2011thesis} \\
Kimbrough et al.\ (2003), two-population GA, GECCO & \nolinkurl{10.1007/3-540-45105-6_123} &
  \citep{kimbrough2008fi2pop} \\
Takahama and Sakai (2006), $\varepsilon$-constrained differential evolution, CEC &
  \nolinkurl{10.1109/CEC.2006.1688283} & \citep{mezuramontes2011survey} \\
Rothberg (2007), evolutionary polishing of MIP solutions, IJOC & \nolinkurl{10.1287/ijoc.1060.0189} &
  \citep{berthold2006primal,mexi2026rexi} \\
Danna et al.\ (2005), relaxation-induced neighbourhoods, Math.\ Prog. & \nolinkurl{10.1007/s10107-004-0518-7} &
  \citep{berthold2006primal,hendel2022alns} \\
French et al.\ (2001), hybrid GA and branch-and-bound, J.\ Heuristics & \nolinkurl{10.1023/A:1011921025322} &
  \citep{puchinger2005combining} \\
Fischetti and Fischetti (2018), Matheuristics, book chapter & \nolinkurl{10.1007/978-3-319-07124-4_14} & --
  \\
Boschetti et al.\ (2009), Matheuristics, HM 2009 & \nolinkurl{10.1007/978-3-642-04918-7_13} & -- \\
Davies et al.\ (2024), ViolationLS: feasibility jump in CP-SAT & \nolinkurl{10.1007/978-3-031-60597-0_16} &
  -- \\
Lodi and Tramontani (2013), MIP performance variability, TutORials & \nolinkurl{10.1287/educ.2013.0112} &
  --$^{b}$ \\
\bottomrule
\multicolumn{3}{p{16.1cm}}{$^{a}$A comparator option in our code; checked in an archived former author page, whose
legal status is open. $^{b}$Reported free at the publisher; not verified.} \\
\end{tabular}
\end{table}


\end{document}
